\chapter{Quadratic Forms over $\Q$ and $\Q_p$}

\section{Quadratic Forms}\label{sec:quadratic_form}

The following definition is available in \mathlib under the name \code{QuadraticForm}:
\begin{definition}\label{def:quadratic_form}\leanok
Let $M$ be a module over a commutative ring $A$. A function $f : M \to A$ is called a \emph{quadratic form} on $M$ if 
\begin{enumerate}
\item $f(ax) = a^2f(x)$ for all $a \in A, x \in M$,
\item The function $(x, y) \mapsto f(x + y) - f(x) - f(y)$ is a bilinear form.
\end{enumerate}
\end{definition}

The \code{QuadraticForm} folder of \mathlib contains definitions and results about quadratic forms, including the bilinear map associated to a quadratic form (see \code{QuadraticMap.associated}) and the definition and characterization of nondegenerate quadratic form (see \code{QuadraticMap.Nondegenerate}).

\begin{definition}\label{def:isotropic}
\lean{QuadraticForm.Isotropic}
\uses{def:quadratic_form}
\leanok
A quadratic form $f : M \to A$ is called \emph{isotropic} if there exists a nonzero $x \in M$ such that $f(x) = 0$.
\end{definition}

\subsection{Orthogonal basis}\label{subsec:ortho_basis}

Let $K$ be a field in which $2$ is invertible and let $V$ be a $K$-vector space.

The definition of orthogonal basis for $V$ is available in Mathlib as \code{LinearMap.IsOrtho}$_i$. The lemma \code{LinearMap.BilinForm.exists_orthogonal_basis} proves that for any symmetric bilinear form $B$ on $V$ there exists an orthogonal basis with respect to $B$.

\begin{definition}\label{def:contiguous}
\lean{Module.Basis.IsContiguous}
\leanok
Two bases of $V$ are called \emph{contiguous} if they have an element in common.
\end{definition}

\begin{theorem}\label{thm:exists_chain}
\lean{Module.Basis.chain_of_nondegenerate}
\uses{Module.Basis.Chain}
\leanok
Let $V$ be a $K$-vector space of dimension at least $3$, let $Q$ be a nondegenerate quadratic form on $V$, and let $b, b'$ be two $Q$-orthogonal bases of $V$. Then there exists a finite sequence $(b^{(0)}, \cdots, b^{(m)})$ of orthogonal bases of $V$ such that $b^{(0)} = b$, $b^{(m)} = b'$ and $b^{(i)}$ is contiguous with $b^{(i + 1)}$ for $0 \le i < m$.
\end{theorem}

\begin{proof}
\uses{Module.Basis.exists_const}
See \cite[\S IV.1.4, Theorem 2]{Ser73}.
\end{proof}

\begin{definition}\label{def:chain}
\lean{Module.Basis.Chain}
\leanok
A sequence of bases $(b^{(0)} = b, \cdots, b^{(m)} = b')$ as in Theorem \ref{thm:exists_chain} is called a \emph{chain} of orthogonal bases contiguously relating $b$ to $b'$.
\end{definition}

\begin{lemma}\label{lem:exists_const}
	\lean{Module.Basis.exists_const}
	\leanok
\end{lemma}

\begin{proof}
	See \cite[\S IV.1.4, Lemma after Theorem 2]{Ser73}.
\end{proof}

\subsection{Translations}\label{subsec:translations}

Let $R$ be a commutative ring, let $V$ and $W$ be $R$-modules and let $Q$ and $Q'$ be quadratic forms on $V$ and $W$, respectively. Then $Q \dotplus Q'$ denotes the quadratic form on $V \times W$ obtained by applying $Q$ on $V$ and $Q'$ on $W$. This is available in Mathlib as \texttt{QuadraticMap.prod}. Similarly, $Q \dotminus Q'$ denotes the form $Q \dotplus (-Q')$ on $V \times W$.

\begin{definition}\label{def:hyperbolic}
\lean{QuadraticForm.IsHyperbolic}
A quadratic form is \emph{hyperbolic} if it is equivalent to the form $X^2 - Y^2$.
\end{definition}

\begin{lemma}\label{lem:XY_hyperbolic}
\lean{QuadraticForm.XY_isHyperbolic}
\uses{def:hyperbolic}
The quadratic form $XY$ is hyperbolic.
\end{lemma}

% NOTE: if needed, assume $R$ is a field.
\begin{definition}\label{def:represents}
\lean{QuadraticForm.represents}
Let $R$ be a commutative ring, let $V$ be an $R$-module and let $Q$ be a quadratic form on $V$. We say that $Q$ \emph{represents} an element $r \in R$ if there exists $x \in V$ such that $Q(v) = r$ and $v \ne 0$.
\end{definition}

\begin{lemma}\label{lem:represents_zero}
\uses{def:represents}
Let $R$ be a commutative ring, let $V$ be an $R$-module and let $Q$ be a quadratic form on $V$. Then $Q$ \emph{represents} $0$ if and only if it is isotropic.
\end{lemma}
\begin{proof}
This follows immediately from the definitions.
\end{proof}

% NOTE: if needed, assume $R$ is a field.
\begin{proposition}\label{prop:sum_hyperbolic}
\lean{QuadraticForm.equivalent_hyperbolic_add}
\uses{def:hyperbolic, def:represents}
Let $R$ be a commutative ring. If a quadratic form $Q$ on an $R$-module $V$ represents $0$ and is nondegenerate, then $Q$ is equivalent to $H + Q'$, where $H$ is hyperbolic. Moreover, $f$ represents all elements of $R$.
\end{proposition}
\begin{proof}
Fix $x \in V$ such that $Q(x) = 0$. Since $Q$ is nondegenerate, there exists $z \in V$ such that $x \cdot z = 1$. The element $y = 2z - (z\cdot z)x$ is isotropic and $x \cdot y = 2$.
The restriction of $Q$ to the submodule $U := Rx \oplus Ry$ is the desired hyperbolic form $H$.
\end{proof}

\begin{corollary}\label{cor:represents_iff}
\lean{QuadraticForm.represents_iff_sub_isotropic}	
Let $K$ be a field, let $r \in K\^times$ and let $Q$ be a quadratic form on a free $K$-vector space $V$. Then $Q$ represents $r$ if and only if the form $Q' := Q \dotminus rZ$ represents $0$.
\end{corollary}
\begin{proof}
\uses{prop:sum_hyperbolic}
The forward implication is clear. For the second one, fix a basis for $V$ and assume that $Q' := Q \dotminus rZ$ has a nontrivial zero $(x_1, \cdots, x_n, z) \in V \times K$. Then either $z = 0$, in which case $Q$ represents zero and thus also $r$, or $Q(x_1/z, \cdots, x_n/z) = r$.
\end{proof}

\begin{corollary}\label{cor:prod_isotropic_iff}
\lean{QuadraticForm.prod_isotropic_iff, QuadraticForm.prod_isotropic_iff}	
Let $K$ be a field and $Q$ and $Q'$ be nondegenerate quadratic forms over $K$ of rank at least $1$. Then the following properties are equivalent:
\begin{enumerate}[a)]
\item $Q \dotminus Q'$ represents $0$.
\item There exists $r \in K^\times$ which is represented by $Q$ and $Q'$.
\item There exists $r \in K^\times$ such that $Q \dotminus rZ^2$ and $Q' \dotminus rZ^2$ represent $0$.
\end{enumerate}
\end{corollary}

\begin{proof}
\uses{cor:represents_iff}
The equivalence (b) $\iff$ (c) follows from Corollary \ref{cor:represents_iff}. The implication (b) $\implies$ (a) is trivial. To show (a) $\implies (b)$, write a nontrivial zero of $Q \dotminus Q'$ in the form $(x, y)$ with $Q(x) = Q'(y)$. If $Q(x) \ne 0$, we are done. If $Q(x) = 0$, then $Q$ represents $0$ and thus by Proposition \ref{prop:sum_hyperbolic} also all elements of $K$. In particular, it represents all nonzero values taken by $Q'$. 
\end{proof}

\section{Quadratic Forms over $\Q_p$}\label{sec:quadratic_form_padic}

\begin{definition}\label{def:discr}
	\uses{def:Hilbert_symbol}
	Let $p$ be a prime number and let $f$ be a quadratic form of rank $n$ over $\Q_p$. Take a diagonal quadratic form $a_1X_1^2 + \cdots + a_nX_n^2$ which is equivalent to $f$; then the \emph{discriminant} of $f$ is the element $d(f) = a_1 \cdots a_n \in \Q_p^\times/{\Q_p^\times}^2$.
	
Note that \mathlib contains \code{QuadraticMap.discr}, as an element of $\Q_p$.
\end{definition}

\begin{definition}\label{def:HM_inv}
\uses{def:Hilbert_symbol}
\lean{QuadraticForm.HasseMinkoskiInvariant}
Let $p$ be a prime number and let $f$ be a quadratic form of rank $n$ over $\Q_p$. Take a diagonal quadratic form $a_1X_1^2 + \cdots + a_nX_n^2$ which is equivalent to $f$; then the \emph{Hasse--Minkowski invariant} of $f$ is the element
\[ \epsilon(f) := \prod_{i < j}(a_i, a_j).\]
\end{definition}

\begin{lemma}\label{lem:HM_inv_well_defined}
\uses{def:HM_inv}
\lean{QuadraticForm.HasseMinkoskiInvariant.eq_of_equivalent}
The value $\epsilon(f)$ does not depend on the choice of diagonal quadratic form equivalent to $f$.
\end{lemma}
\begin{proof}
Let $b = (b_1, \cdots, b_n)$ be an orthogonal basis for $V$ such that $Q(x) = a_1 x_1^2 + \cdots + a_n x_n^2$ for $x = \sum x_i b_i$. We denote  $\epsilon(b) := \prod_{i < j}(a_i, a_j)$. 

If $n=1$, then $\epsilon(f) = 1$. If $n=2$, then $\epsilon(f) = 1$ if and only if the form $Z^2 - a_1 X^2 - a_2 Y^2$ represents $0$. By {\color{red}{TODO}}, this is equivalent to $a_1 X^2 + a_2 Y^2$ representing $1$, which means that there exists $v \in V$ with $Q(v) = 1$; this is independent on $b$. 
For $n \ge 3$ we use induction on $n$. By \ref{thm:exists_chain}, it suffices to prove that for any pair of contiguous bases $b$ and $b'$, we have $\epsilon(b) = \epsilon(b')$. 
Thanks to the symmetry of the Hilbert symbol, $\epsilon(b)$ is unchanged by permutation of the basis elements, so we may assume that $b_1 = b_1'$.
Then $b_1' \cdot b_1' = a_1$, and we can write 
\begin{align*}
	\epsilon(b) &= (a_1, a_2 \cdots a_n) \prod_{2 \le i < j} (a_i, a_j) = \\
	&= (a_1, d(Q) a_1) \prod_{2 \le i < j} (a_i, a_j).
\end{align*}
Analogously, \[\epsilon(b') = (a_1', d(Q) a_1) \prod_{2 \le i < j} (a_i', a_j'). \]
The result follows by combining $a_1 = a_1'$ and the induction hypothesis applied to the orthogonal complement of $e_1$.
\end{proof}

\begin{lemma}\label{lem:HM_inv_pm_one}
\uses{def:HM_inv}
\lean{QuadraticForm.HasseMinkoskiInvariant.eq_one_or_neg_one}
For every quadratic form $f$, we have $\epsilon(f) = \pm 1$.
\end{lemma}
\begin{proof}
Each term in the product is equal to $\pm 1$.
\end{proof}

\section{Hasse--Minkowski}\label{sec:Hasse--Minkowski}

The goal of this section is to prove the Hasse--Minkowski theorem over the rational numbers. To simplify the formalization process, rather than including a single long proof of this theorem, we will split it into smaller lemmas that can be formalized independently.

Denote by $V$ the union of the set of prime numbers and the symbol $\infty$, and denote $\Q_\infty = \R$. 

Let $f$ be a quadratic form over $\Q$. For each $v \in V$, the injection $\Q \to \Q_v$ allows one to view $f$ as a quadratic form over $\Q_v$, which we will denote $f_v$.

\begin{definition}\label{def:everywhere_locally_isotropic}
\lean{QuadraticForm.EverywhereLocallyIsotropic}
\uses{def:isotropic}
\leanok
A quadratic form $f$ with coefficients in $\Q$ is \emph{everywhere locally isotropic} if $f_v$ is isotropic for all $v \in V$.
\end{definition}

\begin{theorem}[Hasse--Minkowski]\label{thm:Hasse--Minkowski}
\lean{QuadraticForm.HasseMinkowski}
\uses{lem:HM_of_degenerate, lem:ev_loc_isotr_of_isotr}
\leanok
Let $f$ be a quadratic form over $\Q$. Then $f$ is isotropic if and only if $f$ is everywhere locally isotropic. In other words, $f$ has a global zero if and only if $f$ has everywhere a local zero. 
\end{theorem}
\begin{proof}
Follows from the results below.
\end{proof}

\begin{theorem}[Hasse--Minkowski]\label{lem:HM_of_degenerate}
	\lean{QuadraticForm.HasseMinkowski}
	\uses{lem:ev_loc_isotr_of_isotr, lem:isotr_of_rk_zero, lem:isotr_of_rk_one, lem:isotr_of_rk_two, lem:isotr_of_rk_three, lem:isotr_of_rk_four, lem:isotr_of_fiv_le_rk}
	\leanok
	Let $f$ be a degenerate quadratic form over $\Q$. Then $f$ is isotropic if and only if $f$ is everywhere locally isotropic. In other words, $f$ has a global zero if and only if $f$ has everywhere a local zero. 
\end{theorem}
\begin{proof}
	
\end{proof}

From now on, given \ref{lem:HM_of_degenerate}, we may assume that $f$ is a nondegenerate quadratic form.

\begin{lemma}\label{lem:ev_loc_isotr_of_isotr}
\lean{QuadraticForm.Isotropic.everywhereLocallyIsotropic}
\uses{def:everywhere_locally_isotropic}
\leanok
Let $f$ be an isotropic quadratic form over $\Q$. Then $f$ is everywhere locally isotropic.
\end{lemma}
\begin{proof}
	
\end{proof}

\begin{lemma}\label{lem:isotr_of_rk_zero}
	\lean{QuadraticForm.EverywhereLocallyIsotropic.isotropic_of_rank_zero}
	\uses{def:everywhere_locally_isotropic}
	\leanok
	Let $f$ be a quadratic form over $\Q$ of rank zero which is everywhere locally isotropic. Then $f$ is isotropic.
\end{lemma}
\begin{proof}
Suffices to see that every rank zero quadratic form over a field is anisotropic.
\end{proof}

\begin{lemma}\label{lem:equiv_sum_squares}
\lean{QuadraticForm.equivalent_weightedSumSquares_units_of_nondegenerate}
\uses{def:quadratic_form}
\leanok
Let $K$ be a field of characteristic not equal to $2$, let $M$ be a $K$-vector space of positive dimension $n$ and let $f : M \to K$ be a
nondegenerate quadratic form. Then $f$ is equivalent to a form $X_1^2 +  a_2X_2^2 + \cdots + a_nX_n^2$, for some $a_i \in \Q^\times$.
\end{lemma}
\begin{proof}
Note that \mathlib contains a similar statement\footnote{\code{QuadraticForm.equivalent_weightedSumSquares_units_of_nondegenerate}}, where the coefficient of $X_1$ is not assumed to be $1$.
\end{proof}

\begin{lemma}\label{lem:isotr_of_rk_one}
	\lean{QuadraticForm.EverywhereLocallyIsotropic.isotropic_of_rank_one}
	\uses{def:everywhere_locally_isotropic}
	\leanok
	Let $f$ be a quadratic form over $\Q$ of rank one which is everywhere locally isotropic. Then $f$ is isotropic.
\end{lemma}
\begin{proof}
Suffices to see that every rank zero quadratic form of rank one over a field is anisotropic.
\end{proof}

\begin{lemma}\label{lem:isotr_of_rk_two}
	\lean{QuadraticForm.EverywhereLocallyIsotropic.isotropic_of_rank_two}
	\uses{def:everywhere_locally_isotropic, lem:equiv_sum_squares}
	\leanok
	Let $f$ be a quadratic form over $\Q$ of rank two which is everywhere locally isotropic. Then $f$ is isotropic.
\end{lemma}
\begin{proof}
$f$ is equivalent to $X_1^2 - aX_2^2$, for some nonzero $a$, which is positive because $f_\infty$ represents $0$. Write
\[ a = \prod_p p^{v_p(a)}.\]
Since $f_p$ represents $0$, $a$ is a square in $\Q_p$, so $v_p(a)$ is even. It follows that $a$ is a square in $\Q$ and $f$ represents $0$.
\end{proof}

\begin{lemma}\label{lem:isotr_of_rk_three}
	\lean{QuadraticForm.EverywhereLocallyIsotropic.isotropic_of_rank_three}
	\uses{def:everywhere_locally_isotropic, lem:equiv_sum_squares}
	\leanok
	Let $f$ be a quadratic form over $\Q$ of rank three which is everywhere locally isotropic. Then $f$ is isotropic.
\end{lemma}
\begin{proof}
	
\end{proof}

\begin{lemma}\label{lem:isotr_of_rk_four}
	\lean{QuadraticForm.EverywhereLocallyIsotropic.isotropic_of_rank_four}
	\uses{def:everywhere_locally_isotropic, lem:equiv_sum_squares}
	\leanok
	Let $f$ be a quadratic form over $\Q$ of rank four which is everywhere locally isotropic. Then $f$ is isotropic.
\end{lemma}
\begin{proof}
	
\end{proof}

\begin{lemma}\label{lem:isotr_of_five_le_rk}
	\lean{QuadraticForm.EverywhereLocallyIsotropic.isotropic_of_five_le_rank}
	\uses{def:everywhere_locally_isotropic, lem:equiv_sum_squares, thm:padic_squares_open, lem:approx_thm, def:HM_inv}
	\leanok
	Let $f$ be a quadratic form over $\Q$ of rank $n \ge 5$ which is everywhere locally isotropic. Then $f$ is isotropic.
\end{lemma}
\begin{proof}
The proof uses induction on $n$. Write $f = h - g$, where $h = a_1X_1^2 + a_2X_2^2$, $g = -(a_3X_3^2 + \cdots + a_nX_n^2)$.
Let $S$ be the finite subset of $V$ consisting of $\infty, 2$, and the primes $p$ such that $v_p(a_i) \ne 0$ for $i \ge 3$.
Let $v \in S$. Since $f_v$ represents $0$, there exists $c_v \in \Q_v^\times$ which is represented in $\Q_v$ by $h$ and $g$, that is, there exist $x_i^v \in \Q_v$, $i=1, \dots, n$ (with $x_i^v \ne 0$ for $i=1, 2$, and some $i \ge 3$) such that 
\[ h(x_1^v, x_2^v) = c_v = g(x_3^v, \cdots, x_n^v).\]
By \ref{thm:padic_squares_open}, the squares of $\Q_p^\times$ form an open set. Combining this with the approximation theorem \ref{lem:approx_thm}, we deduce the existence of $x_1, x_2 \in \Q^\times$ such that, denoting $c := h(x_1, x_2)$, one has $c/c_v \in {\Q_v^\times}^2$ for all $v \in S$.

Consider the form $f_1 := cZ^2 - g$. For each $v \in S$, $g$ represents $c_v$ in $\Q_v$, and hence also $c$ because $c/c_v \in {\Q_v^\times}^2$, hence $f_1$ represents $0$ in $\Q_v$. For $v \notin S$, the coefficients $-a_3, \dots, -a_n$ of $g$ are $v$-adic units, so the discriminant $d_v(g)$ is a $v$-adic unit and $\epsilon_v(g)=1$.


In all cases, $f_q$ represents $0$ in all $\Q_v$, and since the rank of $f$ is $n-1$, the inductive hypothesis shows that $f_1$ represents $0$ in $\Q$, or equivalently, $g$ represents $c$ in $\Q$. But $h$ also represents $c$, so $f$ represents $0$.

\end{proof}